%-------------------------
% Resume in Latex
% Author : Audric Serador
% Inspired by: https://github.com/sb2nov/resume
% License : MIT
%------------------------
% GENERATED FILE - do not edit by hand.
% Produced by resume_builder from a render-ready resume JSON.

%----------FONT TABLE----------%
% Master font table - modify these to change fonts throughout the resume.
%
% +---------------------+--------------+-------------------------------+
% | Element             | Command      | Default                       |
% +---------------------+--------------+-------------------------------+
% | Font family (NFSS)  | \resumefam   | phv  (Helvetica)              |
% | Heading shape       | \resumesc    | \scshape (small caps)         |
% | Name size           | \resumename  | \Large  (14.4pt)              |
% | Section size        | \resumesec   | \Large  (14.4pt)              |
% | Sub/item size       | \resumesub   | \small  (10pt)                |
% | Base body size      | (fixed)      | 10pt  (in \documentclass)     |
% +---------------------+--------------+-------------------------------+
%
% NFSS font family codes (each needs its \usepackage loaded below):
%   phv   : Helvetica (sans-serif, modern/tech - default)
%   ptm   : Times       (serif, traditional/academic)
%   lmr   : Latin Modern Roman (LaTeX default serif)
%   lmss  : Latin Modern Sans  (LaTeX default sans)
%
% To switch font family:
%   1. Change \resumefam below to the desired code.
%   2. In the "Apply font family" block further down, swap the \usepackage
%      line (e.g. \usepackage{helvet} -> \usepackage{mathptmx}).

\newcommand{\resumefam}{phv}
\newcommand{\resumesc}{\scshape}
\newcommand{\resumename}{\Large\resumesc}
\newcommand{\resumesec}{\Large\resumesc}
\newcommand{\resumesub}{\small}

\documentclass[letterpaper,10pt]{article}

\usepackage{fontawesome5}
\usepackage{latexsym}
\usepackage[left=0.75in, right=0.75in, top=0.5in, bottom=0.5in]{geometry}
\usepackage{titlesec}
\usepackage{marvosym}
\usepackage[usenames,dvipsnames]{color}
\usepackage{verbatim}
\usepackage{enumitem}
\usepackage[hidelinks]{hyperref}
\usepackage{fancyhdr}
\usepackage[english]{babel}
\usepackage{tabularx}
\usepackage{setspace}

% pdfTeX-only ATS support (skipped under XeTeX/tectonic)
\ifdefined\pdfgentounicode
  \input{glyphtounicode}
  \pdfgentounicode=1
\fi

% Apply font family: professional sans-serif (Helvetica)
% To switch fonts, swap the \usepackage and the family code in \resumefam above.
\usepackage{helvet}
\renewcommand{\familydefault}{\resumefam}

\setstretch{1}

\pagestyle{fancy}
\fancyhf{} % clear all header and footer fields
\fancyfoot{}
\renewcommand{\headrulewidth}{0pt}
\renewcommand{\footrulewidth}{0pt}

\urlstyle{same}

\raggedbottom
\raggedright
\setlength{\tabcolsep}{0in}

% Sections formatting
\titleformat{\section}{
  \vspace{-4pt}\resumesec\raggedright
}{}{0em}{}[\color{black}\titlerule\vspace{-5pt}]

% Restore filled-circle bullets (Helvetica defaults to en-dash)
\renewcommand{\labelitemi}{$\bullet$}
\renewcommand{\labelitemii}{$\bullet$}

%-------------------------%
\begin{document}

%-------------------------------------------%
%%%%%%  RESUME STARTS HERE  %%%%%

%----------HEADING----------%
\begin{center}
    \textbf{\resumename Abhay Harish Kashyap} \\ \vspace{1pt}
    \href{mailto:abhayhk2001@gmail.com}{\underline{abhayhk2001@gmail.com}} \quad
    +1 217 305 1018 \quad
    \href{https://www.linkedin.com/in/abhay-h-kashyap/}{\underline{Linkedin}} \quad
    \href{https://github.com/abhayhk2001}{\underline{Github}}
\end{center}


%-----------EDUCATION-----------%
\section{Education}
    \begin{itemize}[leftmargin=0.15in, label={}]
    \vspace{-2pt}\item\begin{tabular*}{\linewidth}[t]{@{\extracolsep{\fill}}lr}
      \textbf{University of Illinois Urbana-Champaign} & Champaign, USA \\
      \textit{\resumesub Master of Computer Science} & \textit{\resumesub August 2026 - December 2027} \\
    \end{tabular*}\vspace{-7pt}

    \vspace{-2pt}\item\begin{tabular*}{\linewidth}[t]{@{\extracolsep{\fill}}lr}
      \textbf{RV College of Engineering} & Bengaluru, India \\
      \textit{\resumesub Bachelor in Computer Science and Engineering (GPA: 9.22/10.0)} & \textit{\resumesub August 2019 - May 2023} \\
    \end{tabular*}\vspace{-7pt}
\end{itemize}

%-----------EXPERIENCE-----------%
\section{Work Experience}
\begin{itemize}[leftmargin=0.15in, label={}]

    \vspace{-2pt}\item\begin{tabular*}{\linewidth}[t]{@{\extracolsep{\fill}}lr}
      \textbf{Qualcomm India Private Limited} & Bengaluru, India \\
      \textit{\resumesub Software Engineer} & \textit{\resumesub January 2023 - July 2026} \\
    \end{tabular*}\vspace{-7pt}
    \begin{itemize}
        \item \resumesub{Designed and delivered a production-ready SoC upgrade framework under tight deadlines and limited testing resources}
        \item \resumesub{Optimized the data-relay mechanism, increasing Bluetooth throughput by 24\% and reducing system load during upgrades}
        \item \resumesub{Refactored core component of DFU module, delivering two major client releases on schedule and reducing integration defects while improving system robustness}
        \item \resumesub{Introduced GenAI-based tools for log-analysis automation which reduced triage time, and creating code documentation}
    \end{itemize}\vspace{-5pt}

    \vspace{-2pt}\item\begin{tabular*}{\linewidth}[t]{@{\extracolsep{\fill}}lr}
      \textbf{Epsilon Pvt Ltd} & Bengaluru, India \\
      \textit{\resumesub Software Engineering Intern} & \textit{\resumesub August 2022 - January 2023} \\
    \end{tabular*}\vspace{-7pt}
    \begin{itemize}
        \item \resumesub{Benchmarked HBase, Apache Phoenix and Solr for large-scale data retrieval and analytics, evaluating filtering, aggregation, and multi-condition queries across varying dataset sizes.}
        \item \resumesub{Optimized database query performance through primary/secondary indexing, query restructuring, and selective column retrieval, reducing representative Phoenix query execution from 19.4s to 0.37s.}
        \item \resumesub{Built and evaluated a distributed data-processing environment using Hadoop/HDFS, HBase, Phoenix, Solr, and Docker, analyzing performance across storage and query architectures and documented benchmarking results for large datasets.}
    \end{itemize}\vspace{-5pt}
\end{itemize}

%-----------PROJECTS-----------%
\section{Project Experience}
\begin{itemize}[leftmargin=0.15in, label={}]
    \item\noindent\resumesub\textbf{Business Documents Data Recognition and Table Extraction using Deep Learning Techniques}\vspace{-7pt}
    \begin{itemize}
        \item \resumesub{Integrated DETR-based table extraction and YoloV5 object detection, achieving mAP values of 84.36\% and 80.92\% across diverse business documents}
        \item \resumesub{Presented at International Conference of Artificial Intelligence and Computing (ICAIC) 2025. \href{https://rdcu.be/wZSePwoXqLVa}{\underline{Paper Link}}}
    \end{itemize}\vspace{-5pt}

    \item\noindent\resumesub\textbf{Novel Biomarker Prediction for Lung Cancer Using Random Forest Classifiers}\vspace{-7pt}
    \begin{itemize}
        \item \resumesub{Evaluated multiple ML and ensemble algorithms on gene expression analysis data to detect NSCLC/SCLC causing genes, and Random Forest predict potential biomarkers. Optimized the model to achieve 0.93 AUC and 87\% accuracy}
        \item \resumesub{Published in Cancer Informatics. \href{https://journals.sagepub.com/doi/10.1177/11769351231167992}{\underline{Paper Link}}}
    \end{itemize}\vspace{-5pt}

    \item\noindent\resumesub\textbf{Conditional Independence Testing using RCoT}\vspace{-7pt}
    \begin{itemize}
        \item \resumesub{Modified RCoT algorithm to run on high-performance computing systems (HPCC), reducing runtime by 50\% using synthetic data}
        \item \resumesub{Published in JAIT Journal. \href{http://www.jait.us/shw-229-1347-1.html}{\underline{Paper Link}}}
    \end{itemize}\vspace{-5pt}
\end{itemize}

%-----------SKILLS-----------%
\section{Skills}
    \begin{itemize}[leftmargin=0.15in, label={}]
	\resumesub{\item{
			\textbf{Programming Languages}: Python, Java, C/C++, SQL, Dart, R, Rust \\
			\textbf{Frameworks}: ReactJS, NodeJS, Kotlin, Tensorflow, OpenSSL, Redis, PostgreSQL, SQLite, MongoDB \\
			\textbf{Tools}: Jira, Git/Github, GCP, AWS, Shell Scripting, Jenkins, Perforce, Kubernetes, Terraform, Ansible \\
			\textbf{Machine Learning}: CNN, RNN, LSTM BERT, Transformers, PyTorch, Keras, TensorFlow, Spacy, Langchain, GAN, VAE, RAG
	}}
    \end{itemize}


%-------------------------------------------%
\end{document}
